\section{Application proof-obligation template}
\label{app:obligations}

A new application should fill the following checklist.  It is designed to keep
mathematical status and compiler status separate.

\subsection{Public theorem}

\begin{enumerate}
  \item State the exact ordinary \(\iPRA\) formula, including free parameters,
        coding conventions, and requested precision.
  \item State whether the endpoint is an open sequent, closed sentence, or a
        backend translation.
  \item Identify every public witness or code that must be named.
\end{enumerate}

\subsection{Representations}

For each mathematical object:
\begin{enumerate}
  \item give its numerical code and primitive-recursive operations;
  \item give validity and equality/compatibility formulas;
  \item specify failed or singular-case conventions;
  \item list checked upstream symbols that are reused;
  \item provide reflection theorems for every atom used in bounded formulas.
\end{enumerate}

\subsection{Plan nodes}

For each node record:
\begin{enumerate}
  \item node identifier and kind;
  \item dependency formulas;
  \item output formula;
  \item local \(\Prf\;\PRAAx\) theorem name or explicit residual status;
  \item certificate scope: none, local, observed, or named;
  \item transfer summary with a human proof of soundness;
  \item imports and replay command.
\end{enumerate}

\subsection{Loops and relations}

For every loop:
\begin{enumerate}
  \item list state registers and step terms;
  \item display the quantifier-free invariant after Boolean compilation;
  \item provide base and step derivations;
  \item provide a fuel or termination-bound theorem.
\end{enumerate}
For every relational cloud:
\begin{enumerate}
  \item state layer bounds and edge bits;
  \item state terminal acceptance;
  \item decide whether the path is local, observed, or named;
  \item state any extension or tail producer separately from finite viability.
\end{enumerate}

\subsection{Quantitative producers}

For every approximation operation:
\begin{enumerate}
  \item state the error predicate;
  \item prove the precision transfer and its monotonicity;
  \item state magnitude, apartness, gap, coverage, and tail requirements;
  \item classify each requirement as checked, to be proved, or residual.
\end{enumerate}

\subsection{Compilation audit}

The generated record should include:
\begin{enumerate}
  \item resource shape, query seed, fixed point, and iteration history;
  \item monotonicity evidence for every transfer;
  \item retained and pruned nodes;
  \item final residual formulas;
  \item exact ordinary endpoint;
  \item theorem names used in the simultaneous-cut schedule;
  \item axiom audit and absence of \code{sorryAx};
  \item backend symbol comparisons and translated endpoint, if any.
\end{enumerate}

\subsection{Status vocabulary}

Use only the following status labels:
\begin{description}
  \item[Checked upstream.] A declaration was built and audited in the pinned
  repository or a recorded compatible baseline.
  \item[Checked in application.] A new declaration was built and audited by
  the application replay.
  \item[Human proof.] A complete nonformal argument is supplied but no Lean
  elaboration is claimed.
  \item[Source candidate.] A no-placeholder source file exists but has not
  been elaborated in the recorded environment.
  \item[Lean blueprint.] Signatures and proof plan exist with explicit
  milestones.
  \item[Executable finite test.] Host code verifies finite examples only.
  \item[Residual.] The exact producer theorem remains an assumption.
\end{description}
This vocabulary prevents a compiled PDF, passing Python test, or no-placeholder
text file from being mistaken for a kernel-checked object theorem.
